\section{Introduction}
\label{sec:intro}

With rapid advancements in robotics and artificial intelligence, robots are moving from industrial settings into everyday human environments~\cite{lecun2022path,gonzalez2021service}, showing promise in areas like daily service, home care, and healthcare. To achieve the vision of human-robot symbiosis, the most challenging issue is how to enable robots to interact and collaborate with humans as naturally and efficiently as possible in daily life.%while ensuring safety and privacy protection.
%The ultimate goal is not just to handle repetitive or dangerous tasks, but to work alongside humans to improve efficiency and quality of life—achieving human-robot symbiosis. 
%Realizing this vision requires harmonious and effective communication and interaction between humans and robots, encompassing both verbal and behavioral exchanges. 
%While large language models~\cite{roumeliotis2023chatgpt,touvron2023llama} have made language-based communication easier, major challenges remain in behavioral interaction~\cite{safavi2024emerging}, which involves empowering robots to perceive and interpret human intentions through observation and reasoning, and to adapt their own actions accordingly—fostering seamless, natural collaboration in shared environments.


To support the development and evaluation of intelligent robotic algorithms, over the past few decades, many robotic simulators~\cite{makoviychuk2021isaac, todorov2012mujoco,igibson, nasiriany2024robocasa} have emerged and been widely used. They are well-suited for tasks such as robotic navigation, manipulation, and control within predefined digital environments containing static objects. However, they struggle to support human-robot interaction (HRI), as modeling human behavior—especially the complex, dynamic, and closed-loop interactive behavior—is extremely challenging. To bridge this gap, recent work~\cite{HumanTHOR, liu2024collabsphere} has shifted toward incorporating real humans into HRI simulations. However, these systems still have limitations: humans generally control virtual avatars through keyboards or VR interfaces to collaborate with simulated robots. The focus is on completing tasks that do not involve close physical interaction. So far, there still lacks 3D physical interaction, preventing truly realistic and immersive human-robot behavioral exchange. %This setup lacks 3D physical interaction, preventing truly realistic and immersive human-robot behavioral exchange.

%Just as humans usually learn and grow through practice, robots are required to gain knowledge and skills through continuous real-world interactions, progressively enhancing their intelligence. Nevertheless, due to concerns about efficiency, cost, and safety, most current research on embodied intelligence avoids direct use of physical robots, relying instead on physical~\cite{makoviychuk2021isaac, todorov2012mujoco} or robotic simulators~\cite{igibson, nasiriany2024robocasa}. These simulators are well-suited for tasks such as navigation, manipulation, and control within predefined digital environments containing static objects. However, they struggle to support human-robot interaction (HRI), as modeling human behavior—especially modeling complex, dynamic, and closed-loop interactive behavior—is extremely challenging.
%Recent efforts~\cite{HumanTHOR, liu2024collabsphere} have begun to address the integration of humans into robotic simulation platforms, these approaches typically depend on keyboard or VR-based control of virtual avatars for task-level collaboration with virtual robots. 
%However, humans cannot interact directly with robots in a physical environment, limiting the emergence of authentic interactive behaviors. More importantly, robots are unable to observe real human actions, making it difficult to develop and evaluate intelligent interactive models for robots.

%To enable seamless human-robot coexistence through natural and trustworthy interaction, it is essential to develop human-in-the-loop interactive platforms. Such system would allow real humans to engage in authentic interactions and provide real-time feedback during interaction with robots, ultimately driving the evolution of robotic interactive capabilities toward true symbiosis.
% \begin{haoran}
To address this gap, our system aims to enable real-time and realistic human–robot interaction through AR. AR allows humans to interact directly with virtual robots, and through repeated interactions, users develop familiarity with robot behaviors. Over time, the system collects human feedback that continuously improves robot motion generation. Together, these elements represent a first step toward achieving human–robot symbiosis.
% \end{haoran}

To facilitate the development, evaluation, and evolution %iteration
of HRI algorithms—ultimately enabling intelligent robots to provide direct services to humans—a simulation system supporting authentic HRI experiences is crucial. However, this presents significant challenges: 
%\textbf{1)} How can we accurately represent realistic human behavior? 
\textbf{1)} How to involve real humans interacting with virtual robots in physical space for authentic interaction experience? \textbf{2)} How to drive virtual robots to respond to human actions in real-time, achieving natural and seamless interactive behaviors? \textbf{3)} How to enable continuous algorithmic self-improvement through ongoing interaction, ultimately creating intelligent robots that meet human expectations?

In this paper, we present SymBridge, a pioneering human-in-the-loop robotic interaction system designed for achieving human-robot symbiosis, as Fig.~\ref{fig:teaser} shows. %First, in SymBridge, we employ a LiDAR-based motion capture system to digitize 3D human motions in real time. This approach not only captures precise movement data but also offers inherent privacy preservation by avoiding the collection of any appearance textures or identifiable visual features. Second, we utilize the augmented reality (AR) technology to connect the real human and virtual robot by rendering robot in physical environment by AR glass to provide authentic interactive experience for users. Third, we propose an efficient and effective robotic interaction model to drive robots naturally responds to human actions for completing diverse interactive behaviors. Unlike other offline motion generation methods~\cite{xu2023interdiff,xu2024regennet}, our interaction model could support real-time, online robot action generation. This capability enables the robot to dynamically adapt to rapid changes in human behavior. Moreover, our method learns hierarchical human behavior features and leverages affordance prediction to guide more precise interaction actions, ultimately achieving a natural and smooth interaction experience.  Finally, SymBridge can automatically collect human feedback throughout the interaction process and use the recorded data to fine-tune the robots' capability, enabling the robot's behaviors to progressively adapt to user preferences, creating a "better with use" robotic evolution mechanism. By this, SymBridge can significantly reduce the costs and risks of developing intelligent interaction models for robots.
%SymBridge integrates three key innovations to address above three challenges. 
SymBridge innovatively addresses the above challenges with the following key characteristics: 

\noindent\textit{\textbf{1) Realistic Interaction Interface}}: SymBridge simulates robot perception with a LiDAR-based motion capture system and presents robot reactions through augmented reality (AR) technology. The LiDAR-based MoCap module implements real-time 3D human motion digitization while robust to varying lighting conditions and ensuring privacy-preserving. The AR interface bridges the physical and virtual realms by rendering virtual robots into the user’s environment via AR glasses, creating an immersive, authentic interactive experience that seamlessly blends human and robotic actions. %In addition, to enable accurate robot perception of 3D human movements, SymBridge leverages a precise LiDAR-based motion capture system for real-time 3D human motion digitization. This solution is robust to varying lighting conditions while ensuring ensures privacy-preserving.

\noindent\textit{\textbf{2) Real-Time Robotic Interaction Model}}:
Unlike offline motion generation methods, SymBridge's interaction model enables online, real-time robot action generation. Its key innovation lies in the synergistic integration of spatial reasoning and temporal understanding. Specifically, it fuses an Affordance Predictor for fine-grained spatial awareness with a Multi-Resolution Human Feature (MRHF) Learner that captures deep temporal context. This holistic approach, distinct from prior work (e.g., JRT~\cite{xu2023joint}, ReGenNet~\cite{xu2024regennet}) that often addresses these aspects separately, allows the model to generate responses that are simultaneously contextually and socially-aware, adapting instantly to rapid changes in human behavior.
% Unlike offline motion generation methods (e.g.,~\cite{xu2023interdiff,xu2024regennet}), SymBridge’s interaction model enables online, real-time robot action generation. By dynamically analyzing hierarchical human behavior features and leveraging affordance prediction, the system generates precise, context-aware robotic responses that adapt instantly to rapid changes in human behavior, ensuring more natural and smooth interactions compared to recent relevant approaches.

\noindent\textit{\textbf{3) Feedback-Driven Adaptation Mechanism}}:
SymBridge continuously collects human feedback during interactions and uses this data to iteratively fine-tune the robot’s capabilities. This creates a "better with use" evolution mechanism, where the robot’s behavior progressively aligns with user preferences, significantly simplifying and boosting the improvement of robotic intelligence.

\noindent Together, these components establish SymBridge as a pioneering system for human-robot interaction, %where privacy, adaptability, and user-centric evolution converge to enable safe, efficient, and intuitive interaction paradigms.
which enables immersive, intuitive, and efficient interaction paradigms.
As a dual bridge between humans and robots and physical and virtual worlds, SymBridge accelerates the realization of human-robot symbiosis, with the potential to pave the way for collaborative ecosystems where robots and humans evolve in tandem.



%In SymBridge, users interact with virtual robots through augmented reality (AR) glasses. As Fig.~\ref{fig:teaser} shows, the platform captures human motions as the robot's perception and presents the virtual robot's reaction to the user from first-person and third-person perspectives. It can automatically collect human feedback throughout the interaction process and use the recorded data to fine-tune the robots' capability, enabling the robot's behaviors to progressively adapt to user preferences, creating a "better with use" robotic system. By this, SymBridge can significantly reduce the costs and risks of developing intelligent interaction models for robots. Simultaneously, users develop a deeper understanding of the robot's intentions and refine their interaction strategies, fostering greater trust and enabling the robot’s seamless integration into daily life. As a dual bridge between humans and robots and physical and virtual worlds, SymBridge accelerates the realization of human-robot symbiosis, with the potential to pave the way for collaborative ecosystems where robots and humans evolve in tandem.

%In particular, to facilitate the validation of SymBridge platform, we provide a novel and effective method for humanoid robot motion generation, enabling the robot to respond promptly to human actions and successfully complete interaction tasks. Unlike other offline motion generation methods~\cite{xu2023interdiff,xu2024regennet}, our robotic interaction model could support real-time, online robot action generation. This capability enables the robot to dynamically adapt to rapid changes in human behavior. Furthermore, our method hierarchically models observed human behavior and leverages affordance prediction to guide more precise interaction actions, ultimately achieving a natural and smooth interaction experience.  


To demonstrate the effectiveness of our solution, we have developed a comprehensive evaluation methodology and benchmark for human-robot interaction, complementing both quantitative and qualitative assessments of the algorithmic components within our system.
Further, we evaluate the efficiency and quality performance of our system when facilitating human users to interact with the virtual robot. 
A user empirical test with 50 participants confirms that human users are satisfied with the realistic and real-time interaction experience via our system. We also deploy the generated virtual robot reactions to a real humanoid robot to examine the realism and fidelity of the robot movements presented with our system. Additionally, by fine-tuning our model with user feedback data, we demonstrate continuous performance improvement, proving our system's capability to enhance robotic interaction models through ongoing interactive data collection. 

%We conduct comprehensive experiments to validate the effectiveness and practicality of the SymBridge platform. First, through both qualitative and quantitative evaluations, we validate the superior performance of our robotic interaction model compared to other alternative solutions.
%Second, we evaluate the time and quality performance of the platform when facilitating human users to interact with the virtual robot. 
%A user empirical test with 50 participants indicates that human users are satisfied with the realistic and real-time interaction experience. We also deploy the virtual robot reactions to a real humanoid robot to examine the realism and fidelity of the interactions presented with our platform. Third, we show the performance improvement of the robotic interaction model after fine-tuning it with more positive and negative interaction data, demonstrating that our platform can assist model's improvement by collecting authentic human-robot interaction data and user feedback. 

%Our contributions are summarized as follows:
%\begin{itemize}
%\item We introduce SymBridge, the first human-in-the-loop cyber-physical system, designed to enable real-time, realistic interactions between human users and virtual robots, and also support feedback-driven adaptation of robotic models.
%\item We propose a novel robotic interaction model, supporting online precise humanoid robot motion generation for versatile human-robot interactions.
%\item We conduct comprehensive evaluations as well as user empirical test and real robot test, to validate the effectiveness and practicality of our system and internal modules.
%\end{itemize}

